\documentclass[11pt]{article}

% --- Encoding / language (XeLaTeX-safe) ---
\usepackage{fontspec}
\usepackage[russian,english]{babel}
\defaultfontfeatures{Ligatures=TeX}
\setmainfont{Times New Roman}

% --- Math / theorem ---
\usepackage{amsmath}
\usepackage{amssymb}
\usepackage{amsthm}
\usepackage{stmaryrd}
\usepackage{breqn}
\usepackage{euscript}
\usepackage{mathrsfs}
\usepackage{enumerate}
\usepackage{empheq}

% --- Graphics / layout ---
\usepackage{graphicx}
\usepackage{wrapfig}
\usepackage{subcaption}
\usepackage{float}
\usepackage{xcolor}
\usepackage{geometry}

\geometry{
    left=1cm, % 2.5
    right=1cm, % 2.5
    top=2cm,
    bottom=2cm,
    bindingoffset=0cm
}

% --- Extra utilities ---
\usepackage{pifont}
\usepackage{lipsum}
\usepackage{todonotes}
\usepackage{csquotes}

% --- Bibliography ---
\usepackage[backend=biber,style=alphabetic]{biblatex}
\addbibresource{references.bib}

% --- Hyperref (must be loaded last) ---
\usepackage{hyperref}

% --- Custom imports ---
\input{notations}

% --- Safe image inclusion (shows empty figure if file missing) ---
\newcommand{\safeincludegraphics}[2][]{%
    \IfFileExists{#2}{%
        \includegraphics[#1]{#2}%
    }{%
        \textcolor{gray}{\rule{0.8\linewidth}{0.6\linewidth}}%
    }%
}

% --- Theorems ---
\newtheorem{proposition}{Proposition}
\newtheorem{theorem}{Theorem}
\theoremstyle{remark}
\newtheorem{remark}{Remark}

% --- Colors ---
\definecolor{darkred}{RGB}{139,0,0}

% --- Date helper ---
\newcommand{\todayYYMMDD}{\the\year-\the\month-\the\day}


\begin{document}

\hfill
% \todayYYMMDD
\begin{center}
% \section*{Custom course or document title}
\textit{Kramers--Smoluchowski limit: from underdamped to overdamped Langevin dynamics}\\
% \textit{Optional subtitle}\\
\textit{Author: jj}\\
% \textit{Supervisor: Name}\\
\textit{Date: \today}
\end{center}
\section{Framework: Langevin dynamics}

Langevin dynamics define a class of second order SDEs characterized by the following (formal) stochastic Newton law:
\[
\gamma \ddot{X}_t = F(t,X_t,\dot{X}_t)\,dt + \sigma(t,X_t,\dot{X}_t)\,\dot{L}_t, \qquad t \ge 0.
\]

From a physical point of view, $X_t=(X_t^r)_{t\ge0}$ models the position of a generic body evolving in $\mathbb{R}^d$,
$\dot{X}_t=(\dot{X}_t^r)_{t\ge0}$ its velocity, $\gamma$ the volume or mass of the body,
$F:[0,\infty)\times\mathbb{R}^d\times\mathbb{R}^d\to\mathbb{R}^d$ an external force field
and $(L_t)_{t\ge0}$ a $m$-dimensional stochastic perturbation weighted by some homogeneous dispersion matrix~$\sigma$.
Setting $V_t=\dot{X}_t$, the above rewrites more rigorously into the $\mathbb{R}^{2d}$-valued SDE system:
\begin{equation}
\begin{cases}
dX_t = V_t\,dt,\\
dV_t = \gamma^{-1}F(t,X_t,V_t)\,dt + \gamma^{-1}\sigma(t,X_t,V_t)\,dL_t,
\end{cases}
\qquad t\ge0.
\tag{1.1}
\end{equation}

In this formulation the coefficient $\gamma^{-1}$ is understood as a friction force.

\begin{remark}
A prototypical instance of \eqref{1.1} is given by the following setting:
\[
L=W \text{ is a $d$-dimensional Brownian motion}, \qquad
F(t,x,v)=-\nabla f(x)-\lambda v,
\]
\[
\sigma(t,x,v)=\sigma I_{d\times d}>0, \quad \sigma \text{ is a scalar constant}.
\tag{1.2}
\]
\end{remark}

In this situation \eqref{1.1} writes as
\begin{equation}
\begin{cases}
dX_t = V_t\,dt,\\
dV_t = -\gamma^{-1}\big(\nabla f(X_t)+V_t\big)\,dt + \gamma^{-1}\sigma\,dW_t,
\end{cases}
\qquad t\ge0.
\tag{1.3}
\end{equation}

Langevin dynamics, originally introduced in~\cite{5}, apply in modern science in a wide range of fields,
including fluid dynamics, biology, economy, social science~\cite{7,6} and, not so long ago, in artificial intelligence~\cite{1}.

\section{Clarification on Pavliotis \cite{8} scaling and limits}

(This section mainly refers to Chapter~6 in~\cite{8}.)  
The Langevin model in~\cite{8} (see Equation~(6.1), p.~181) corresponds to a second order stochastic system describing the motion of a ``particle'' with
\[
d\dot{q}_t = -\nabla U(q_t)\,dt - \eta \dot{q}_t\,dt + \sqrt{\frac{2\eta}{\beta}}\,dW_t,
\qquad t\ge0,
\]
where $q=(q_t)_t$ is the position dynamics of the particle and $\dot{q}=(\dot{q}_t)_t$ the related velocity dynamics.
Motions are assumed to start at the position $q_0$ and velocity $\dot{q}_0$.
(For notation consistency, the notation of the potential force and the friction coefficient are denoted by $\nabla U$ and $\eta$
instead of $\nabla V$ and $\gamma$ in~\cite{8}.)

To this model, we associate the family of rescaled dynamics
\[
q_t^\eta := \lambda_\eta q_{t/\mu_\eta}, \qquad t\ge0.
\]

As such, the corresponding velocity $\dot{q}^\eta$ is given by
$\dot{q}_t^\eta = \frac{\lambda_\eta}{\mu_\eta}\dot{q}_{t/\mu_\eta}$ and so
\[
\begin{aligned}
\dot{q}_t^\eta
&= \frac{\lambda_\eta}{\mu_\eta}\dot{q}_0
 - \frac{\lambda_\eta}{\mu_\eta}\int_0^{t/\mu_\eta} \nabla U(q_s)\,ds
 - \frac{\lambda_\eta}{\mu_\eta}\int_0^{t/\mu_\eta} \dot{q}_s\,ds
 + \frac{\lambda_\eta}{\mu_\eta}\sqrt{\frac{2\eta}{\beta}}\,W_{t/\mu_\eta} \\
&= \frac{\lambda_\eta}{\mu_\eta}\dot{q}_0
 - \frac{\lambda_\eta}{(\mu_\eta)^2}\int_0^{t} \nabla U(q_{s/\mu_\eta})\,ds
 - \frac{\lambda_\eta}{(\mu_\eta)^2}\int_0^{t} \dot{q}_{s/\mu_\eta}\,ds
 + \frac{\lambda_\eta}{\mu_\eta}\sqrt{\frac{2\eta}{\beta}}\,W_{t/\mu_\eta}.
\end{aligned}
\]

\todo[inline]{1 end}

Since (see e.g.~\cite{4}) $(W_{t/\mu_\eta})_{t\ge0}\overset{\mathcal{D}}{=}(W_t/\sqrt{\mu_\eta})_{t\ge0}$, $\dot q^\eta$ can be viewed as solution to the SDE:
\[
d\dot q_t^\eta
=
-\frac{\lambda_\eta}{(\mu_\eta)^2}\nabla U\!\left(q_t^\eta/\lambda_\eta\right)dt
-\frac{\eta}{\mu_\eta}\dot q_t^\eta\,dt
+\frac{\lambda_\eta}{(\mu_\eta)^{3/2}}\sqrt{\frac{2\eta}{\beta}}\,dW_t,
\qquad t\ge0,
\qquad
\dot q_0^\eta=\frac{\lambda_\eta}{\mu_\eta}\dot q_0.
\]

In particular, taking $\lambda_\eta=1$ and $\mu_\eta=1/\eta$, the above writes as
\[
d\dot q_t^\eta
=
-\eta^2\nabla U(q_t^\eta)\,dt-\eta^2\dot q_t^\eta\,dt+\eta^2\sqrt{\frac{2}{\beta}}\,dW_t,
\qquad t\ge0,
\qquad
\dot q_0^\eta=\eta\dot q_0.
\]

Comparing the above with (1.1) (and as intuited above), the (friction) coefficient $\eta^2$ corresponds to the inverse of the mass $\gamma$.
\textcolor{red}{(While the rescaling in \cite{8} are seemingly coherent with our initial Langevin framework, the remaining impact on the initial velocity seems to have been overlooked: the Kramers--Smoluchowski limit $\gamma\downarrow 0$ or equivalently $\eta\uparrow\infty$ makes $\dot q_0^\eta=\eta\dot q_0$ blows up which greatly contrasts with Proposition~1 below \dots\ The initial velocity $\dot q_0$ should have been redefined as $\dot q_0=V_0/\eta$ for some fixed random $V_0\in\mathbb{R}^d$).}

\section{Quantitative Kramers--Smoluchowski limit}

In the proposition below, under the prototypical condition (1.2), we establish the so-called Kramers--Smoluchowski limit asserting the $\gamma\downarrow 0$ (zero-mass) limit of (1.1).

\begin{proposition}
Let $(X^\gamma,V^\gamma)$ be the Langevin dynamic (1.3).
Assume that the initial condition $(X_0^\gamma,V_0^\gamma)$ are independent of $\gamma$, set at $(X_0,V_0)$ with $\mathbb{E}\!\left[|X_0|^2+|V_0|^2\right]<\infty$, and that $\sigma\neq 0$ and $\nabla f$ is Lipschitz-continuous.
Then for all finite time-horizon $T$, $\gamma\le 1$ and integer $p\ge1$,
\begin{equation}
\max_{t\in[0,T]}\mathbb{E}\!\left[|X_t^\gamma-X_t^0|^{2p}\right]=\mathcal{O}(\gamma^p),
\tag{3.4}
\end{equation}
where $X^0$ is the solution to
\begin{equation}
dX_t^0=-\nabla f(X_t^0)\,dt+\sigma\,dW_t,\qquad X_0^0=X_0.
\tag{3.5}
\end{equation}
Moreover,
\begin{equation}
\mathbb{E}\!\left[\max_{t\in[0,T]}|X_t^\gamma-X_t^0|^2\right]=\mathcal{O}(\gamma^{1/2}).
\tag{3.6}
\end{equation}
\end{proposition}

While the Kramers--Smoluchowski limit is more directly linked to (3.6), the proposition is purposely written to emphasize the successive ``easy'' part ((3.4))\footnote{Stated without proof details in \cite{8}, bottom of p.~206.} and more ``tricky'' part (3.6) of the Kramers--Smoluchowski limit.
In particular, for $p=1$, (3.4) is optimal but the discrepancy between the rate $\mathcal{O}(\gamma)$ in (3.4) and the $\mathcal{O}(\gamma^{1/2})$ in (3.6), although reflecting the cost between time marginal and max-in-time mean-square distance (which may be significant), is rather non-intuitive.
The demonstration of (3.6) might be so subject to some proof refinements.

\textcolor{blue}{\textbf{To do (09/10):}
\begin{enumerate}
\item The estimate (3.6) should be extended to larger moment: it is expected that a more general asymptotic of the form ``$\mathbb{E}(\max_{t\in[0,T]}|X_t^\gamma-X_t^0|^{2p})=\mathcal{O}(\gamma^{p/2})$'' holds.
\item Extensions to non-Lipschitz force field should be drawn from the proof details for (3.4) and (3.6) below. Especially given the objective of applying Kramers--Smoluchowski limit to gradient descent methods, the prototypical case where $f$ is a $C^1$ convex function or $\nabla f$ is double-wells potential must be explored.
\end{enumerate}}

\begin{remark}
The discrepancy between the rate $\mathcal{O}(\gamma)$ in (3.4) and the $\mathcal{O}(\gamma^{1/2})$ in (3.6) is rather non-intuitive (and is possibly subject from some refinements) but reflect the cost between time marginal and max-in-time mean-square distance.
\end{remark}

\paragraph{Proof.}
As a preliminary, recall that, under the initial and Lipschitz conditions of the proposition, each Equations (1.1) and (3.5) admit a unique solution with finite second moments at all times.
In particular, as $\nabla f$ is Lipschitz, $\nabla f$ grows (at most) linearly with
\begin{equation}
|\nabla f(x)| \le |\nabla f(0)| + |\nabla f|_{\mathrm{Lip}}\,|x|
\le \max\!\big(|\nabla f(0)|,|\nabla f|_{\mathrm{Lip}}\big)(1+|x|),
\tag{3.7}
\end{equation}

\todo[inline]{2 end}

for $|\nabla f|_{\mathrm{Lip}}=\sup_{x\neq y}\dfrac{|\nabla f(x)-\nabla f(y)|}{|x-y|}$. Consequently (after some calculations on the SDE (3.5), see e.g.~\cite{4}), we get, for all integer $p\ge1$,
\begin{equation}
\max_{t\in[0,T]}\mathbb{E}\!\left[|X_t^0|^{2p}\right]<\infty
\tag{3.8}
\end{equation}

\underline{Proof of (3.4):}
The limit (3.4) is essentially established by stochastic calculus. Define $U_t=\exp\{\gamma^{-1}t\}V_t^\gamma$ and applying It\^o's formula immediately observe that, by It\^o's integration by parts (see Appendix)
\[
dU_t=-\gamma^{-1}\exp\{\gamma^{-1}t\}\nabla f(X_t^\gamma)\,dt+\sigma\gamma^{-1}\exp\{\gamma^{-1}t\}\,dW_t,
\]
or equivalently
\[
U_t=V_0+\gamma^{-1}\int_0^t \exp\{\gamma^{-1}s\}\nabla f(X_s^\gamma)\,ds+\sigma\gamma^{-1}\int_0^t \exp\{\gamma^{-1}s\}\,dW_s.
\]

It thus follows that
\begin{equation}
\begin{aligned}
V_t^\gamma &= \exp\{-\gamma^{-1}t\}U_t \\
&= \exp\{-\gamma^{-1}t\}V_0+\gamma^{-1}\int_0^t \exp\{\gamma^{-1}(s-t)\}\nabla f(X_s^\gamma)\,ds+\sigma\gamma^{-1}\int_0^t \exp\{\gamma^{-1}(s-t)\}\,dW_s,
\end{aligned}
\tag{3.9}
\end{equation}
and, further,
\[
\begin{aligned}
X_t^\gamma &= X_0^\gamma+\int_0^t V_s^\gamma\,ds \\
&= X_0^\gamma+\int_0^t \exp\{-\gamma^{-1}s\}\,ds\,V_0 \\
&\quad -\gamma^{-1}\int_0^t \left(\int_s^t \exp\{\gamma^{-1}(s-r)\}\,dr\right)\nabla f(X_s^\gamma)\,ds+\sigma\gamma^{-1}\int_0^t \left(\int_s^t \exp\{\gamma^{-1}(s-r)\}\,dr\right)\,dW_s \\
&= X_0^\gamma+\gamma V_0\bigl(1-\exp\{-\gamma^{-1}t\}\bigr) \\
&\quad +\int_0^t \bigl(1-\exp\{\gamma^{-1}(s-t)\}\bigr)\nabla f(X_s^\gamma)\,ds+\sigma\int_0^t \bigl(1-\exp\{\gamma^{-1}(s-t)\}\bigr)\,dW_s.
\end{aligned}
\]

Comparing the paths of $X$ and $X^0$, we get
\begin{equation}
\begin{aligned}
X_t^\gamma-X_t^0
&= \gamma V_0\bigl(1-\exp\{-\gamma^{-1}t\}\bigr)+\int_0^t \nabla f(X_s^\gamma)-\nabla f(X_s^0)\,ds \\
&\quad -\int_0^t \exp\{\gamma^{-1}(s-t)\}\nabla f(X_s^\gamma)\,ds-\sigma\int_0^t \exp\{\gamma^{-1}(s-t)\}\,dW_s,
\end{aligned}
\tag{3.10}
\end{equation}
and, equivalently,
\[
X_t^\gamma-X_t^0=\int_0^t \nabla f(X_s^\gamma)-\nabla f(X_s^0)\,ds+\gamma(V_0-V_t^\gamma).
\]

Taking the Euclidean norm $|\cdot|$, applying the triangular inequality, elevating the resulting expression to the power $2p$ and taking the expectation yields
\begin{equation}
\mathbb{E}\!\left[|X_t^\gamma-X_t^0|^{2p}\right]\le 2^{2p-1}\left(\mathbb{E}\!\left[\left|\int_0^t \nabla f(X_s^\gamma)-\nabla f(X_s^0)\,ds\right|^{2p}+\bigl(\gamma|V_0-V_t^\gamma|\bigr)^{2p}\right]\right)
\tag{3.11}
\end{equation}

Using H\"older's inequality and the Lipschitz condition of $\nabla f$,
\[
\mathbb{E}\!\left[\left|\int_0^t \nabla f(X_s^\gamma)-\nabla f(X_s^0)\,ds\right|^{2p}\right]\le |\nabla f|_{\mathrm{Lip}}^{2p}\,t^{2p-1}\int_0^t \left|X_s^\gamma-X_s^0\right|^{2p}\,ds
\]

Meanwhile, recalling (3.9), we have
\[
\begin{aligned}
\mathbb{E}\!\left[\gamma^{2p}|V_0-V_t^\gamma|^{2p}\right]
&\le 3^{2p-1}\Bigl(\gamma^{2p}\bigl(1-\exp\{-\gamma^{-1}t\}\bigr)^{2p}\mathbb{E}\!\left[|V_0|^{2p}\right]+\mathbb{E}\!\left[\left|\int_0^t \exp\{\gamma^{-1}(s-t)\}\nabla f(X_s^\gamma)\,ds\right|^{2p}\right] \\
&\qquad\qquad +\sigma^{2p}\mathbb{E}\!\left[\left|\int_0^t \exp\{\gamma^{-1}(s-t)\}\,dW_s\right|^{2p}\right]\Bigr)
\end{aligned}
\]

\todo[inline]{3 end}

Adding and subtracting $\nabla f(X_s^0)$, using the Lipschitz property of $\nabla f$ and the Jensen inequality, and observing that $\exp\{\gamma^{-1}(s-t)\}\le 1$ for all $s<t$ yield
\[
\mathbb{E}\!\left[\left|\int_0^t \exp\{\gamma^{-1}(s-t)\}\nabla f(X_s^\gamma)\,ds\right|^{2p}\right]
\]
\[
=\mathbb{E}\!\left[\left|\int_0^t \exp\{\gamma^{-1}(s-t)\}\big(\nabla f(X_s^\gamma)-\nabla f(X_s^0)\big)\,ds+\int_0^t \exp\{\gamma^{-1}(s-t)\}\nabla f(X_s^0)\,ds\right|^{2p}\right]
\]
\[
\le 2^{2p-1}|\nabla f|_{\mathrm{Lip}}^{2p}\mathbb{E}\!\left[\int_0^t |X_s^\gamma-X_s^0|^{2p}\,ds\right]+2^{2p-1}\mathbb{E}\!\left[\int_0^t \exp\{2p\gamma^{-1}(s-t)\}|\nabla f(X_s^0)|^{2p}\,ds\right].
\]

Furthermore, according to the property (3.7), we have
\[
\mathbb{E}\!\left[\int_0^t \exp\{\gamma^{-1}(s-t)\}|\nabla f(X_s^0)|^{2p}\,ds\right]
\le |\nabla f(x)| \le \frac{\gamma}{2p}\max\bigl(|\nabla f(0)|^{2p},|\nabla f|_{\mathrm{Lip}}^{2p}\bigr)\bigl(1-\exp\{-2p\gamma^{-1}t\}\bigr)\mathbb{E}\!\left[1+\max_{0\le s\le t}|X_t^0|^{2p}\right],
\]

Finally, using Burkh\"older-Davis-Gundy inequality, we have
\[
\mathbb{E}\!\left[\left|\int_0^t \exp\{\gamma^{-1}(s-t)\}\,dW_s\right|^{2p}\right]
\le (2dp)^p\mathbb{E}\!\left[\left(\int_0^t \exp\{2\gamma^{-1}(s-t)\}\,ds\right)^p\right]
=2^{p-1}(dp)^p\gamma^p\bigl(1-\exp\{-2\gamma^{-1}t\}\bigr)^p.
\]

Combining all the above, we first deduce that
\[
\mathbb{E}\!\left[\gamma^{2p}|V_0-V_t^\gamma|^{2p}\right]\le C\left(\gamma^{2p}\mathbb{E}\!\left[|V_0|^{2p}\right]+\int_0^t |X_s^\gamma-X_s^0|^{2p}\,ds+\gamma\left(\mathbb{E}\!\left[1+\max_{0\le s\le t}|X_s^0|^{2p}\right]+\sigma^{2p}\right)\right)
\]
for some constant $C$ depending only on $p,d$ and $\nabla f$, and coming back to (3.11) - and up to a slight change of the factor $C$, we obtain that, for any $\gamma\le 1$,
\begin{equation}
\mathbb{E}\!\left[|X_t^\gamma-X_t^0|^{2p}\right]\le C\gamma^p+\int_0^t \mathbb{E}\!\left[|X_s^\gamma-X_s^0|^{2p}\right]\,ds.
\tag{3.12}
\end{equation}

Using Gronwall's inequality, (3.4) follows.

\underline{Proof of (3.6):}
Before establishing (3.6), let us point out that, contrary to some papers (e.g. [2]), the limit is so directly deduced from martingale inequalities and needs a few careful steps. Coming back to (??), we deduce that, for all intermediate time-horizon $T_0\in(0,T]$
\[
\max_{t\in[0,T_0]}|X_t^\gamma-X_t^0|^2
\le \gamma^2\max_{t\in[0,T_0]}\bigl(1-\exp\{-\gamma^{-1}t\}\bigr)^2|V_0|^2
+2\max_{t\in[0,T_0]}\int_0^t |X_s^\gamma-X_s^0|^2\,ds
\]
\[
-2\sigma\min_{t\in[0,T_0]}\int_0^t \exp\{\gamma^{-1}(s-t)\}\bigl(X_s^\gamma-X_s^0\bigr)\cdot dW_s
\]
\[
+2\max_{t\in[0,T_0]}\int_0^t \exp\{\gamma^{-1}(s-t)\}\max\bigl(|\nabla f(0)|^2,|\nabla f|_{\mathrm{Lip}}^2\bigr)(1+|X_s^0|)^2\,ds
\]
\[
+\max_{t\in[0,T_0]}\frac{\gamma\sigma^2}{2}\bigl(1-\exp\{-2\gamma^{-1}t\}\bigr)
\]
\[
\le \beta_{T_0}^\gamma+2\int_0^{T_0}|X_s^\gamma-X_s^0|^2\,ds-2\sigma\min_{t\in[0,T_0]}\int_0^t \exp\{\gamma^{-1}(s-t)\}\bigl(X_s^\gamma-X_s^0\bigr)\cdot dW_s
\]
for
\[
\beta_t^\gamma:=\gamma^2|V_0|^2+4\gamma\max\bigl(|\nabla f(0)|^2,|\nabla f|_{\mathrm{Lip}}^2\bigr)\int_0^t (1+|X_s^0|)^2\,ds+\frac{\gamma\sigma^2}{2}.
\]

Taking the expectation of the last inequality gives
\[
\mathbb{E}\!\left[\max_{t\in[0,T_0]}|X_t^\gamma-X_t^0|^2\right]\le \mathbb{E}\!\left[\beta_{T_0}^\gamma\right]+2\int_0^{T_0}\mathbb{E}\!\left[|X_s^\gamma-X_s^0|^2\right]\,ds-2\sigma\mathbb{E}\!\left[\min_{t\in[0,T_0]}\int_0^t \exp\{\gamma^{-1}(s-t)\}\bigl(X_s^\gamma-X_s^0\bigr)\cdot dW_s\right].
\]

At this stage, a simple upper-bound of order $\gamma$ can be derived for $\mathbb{E}[\beta_{T_0}^\gamma]$ (see below) and the main technical lies with the handling of the remaining "stochastic integral"\footnote{The exact terminology should rather be stochastic convolution.}
\[
I_t^\gamma:=\int_0^t \exp\{\gamma^{-1}(s-t)\}\bigl(X_s^\gamma-X_s^0\bigr)\cdot dW_s,\qquad t\ge 0.
\]

\todo[inline]{4 end}

Introducing the process $\bigl(L_t^\gamma:=\int_0^t (X_s^\gamma-X_s^0)\cdot dW_s\bigr)_{t\ge0}$, $I^\gamma$ writes as
\[
I_t^\gamma=\exp\{-\gamma^{-1}t\}L_t^\gamma,\qquad t\ge0,
\]
and it can be observed (by It\^o's integration by parts) that $I^\gamma$ is also solution to the SDE:
\[
dI_t^\gamma=-\gamma^{-1}I_t^\gamma\,dt+dL_t^\gamma,\qquad t\ge0,\qquad I_0=0.
\]

By It\^o's formula, this specifically yields
\[
d|I_t^\gamma|^2=2I_t^\gamma\,dI_t^\gamma+d\langle I^\gamma\rangle_t
=\bigl(|X_t^\gamma-X_t^0|^2-2\gamma^{-1}|I_t^\gamma|^2\bigr)\,dt+2I_t^\gamma\,dL_t^\gamma,
\]
and further
\[
|I_t^\gamma|^2+2\gamma^{-1}\int_0^t |I_t^\gamma|^2\,ds
=\int_0^t |X_t^\gamma-X_t^0|^2\,ds+2\int_0^t I_s^\gamma\,dL_s^\gamma.
\]

Therefore, for all $T_0\le T$,
\[
\mathbb{E}\!\left[\max_{t\in[0,T_0]}|I_t^\gamma|^2\right]
\le \int_0^{T_0}\mathbb{E}\!\left[|X_t^\gamma-X_t^0|^2\right]\,dt
+2\mathbb{E}\!\left[\max_{t\in[0,T_0]}\int_0^t I_s^\gamma\,dL_s^\gamma\right]
\]
where, applying successively the Cauchy-Schwarz and Burkh\"older-Davis-Gundy inequality,
\begin{equation}
\mathbb{E}\!\left[\max_{t\in[0,T_0]}\int_0^t I_s^\gamma\,dL_s^\gamma\right]
\le \left(\mathbb{E}\!\left[\max_{t\in[0,T_0]}\left(\int_0^t I_s^\gamma\,dL_s^\gamma\right)^2\right]\right)^{1/2}
\le 2\left(\mathbb{E}\!\left[\int_0^{T_0}|I_s^\gamma|^2|X_s^\gamma-X_s^0|^2\,ds\right]\right)^{1/2}
\tag{3.13}
\end{equation}

According to Jensen's inequality, $\mathbb{E}[\max_{t\in[0,T_0]}|I_t^\gamma|^2]\ge \bigl(\mathbb{E}[\max_{t\in[0,T_0]}|I_t^\gamma|]\bigr)^2$ and so
\[
\left(\mathbb{E}\!\left[\max_{t\in[0,T_0]}|I_t^\gamma|\right]\right)^2
\le \int_0^{T_0}\mathbb{E}\!\left[|X_t^\gamma-X_t^0|^2\right]\,dt
+2\left(\mathbb{E}\!\left[\int_0^{T_0}|I_s^\gamma|^2|X_s^\gamma-X_s^0|^2\,ds\right]\right)^{1/2}
\]
\[
\le \int_0^{T_0}\mathbb{E}\!\left[|X_t^\gamma-X_t^0|^2\right]\,dt
+\sqrt{2}\left(\int_0^{T_0}\mathbb{E}\!\left[\max_{v\in[0,s]}|I_v^\gamma|^4\right]\,ds
+\int_0^{T_0}\mathbb{E}\!\left[|X_s^\gamma-X_s^0|^4\right]\,ds\right)^{1/2}
\]
using the inequality: for all $a,b\in\mathbb{R}$, $ab\le 1/2(a^2+b^2)$. and taking the square of (3.13), we deduce that
$f(t):=\bigl(\mathbb{E}[\max_{s\in[0,t]}|I_s^\gamma|]\bigr)^4$ satisfies an inequality of the form: for all $0\le T_0\le T$
\[
f(T_0)\le J_\gamma(T_0)+4\int_0^{T_0}f(s)\,ds
\]
for $J_\gamma(t)=4\left(\left(\int_0^{T_0}\mathbb{E}\!\left[|X_s^\gamma-X_s^0|^2\right]\,ds\right)^2+\int_0^{T_0}\mathbb{E}\!\left[|X_s^\gamma-X_s^0|^4\right]\,ds\right)$. Using Gronwall's inequality, it follows
that $f(T_0)\le J_\gamma(T_0)\exp\{4T_0\}$ and further
\[
\mathbb{E}\!\left[\max_{t\in[0,T_0]}|I_t^\gamma|\right]
\le \bigl(J_\gamma(T_0)\bigr)^{1/4}\exp\{T\}
\le \sqrt{2}\left(\int_0^{T_0}\mathbb{E}\!\left[|X_s^\gamma-X_s^0|^4\right]\,ds\right)^{1/4}\exp\{T_0\}.
\]

Coming back to the initial estimate on $\mathbb{E}[\max_{t\in[0,T_0]}|X_t^\gamma-X_t^0|^2]$, we get
\[
\mathbb{E}\!\left[\max_{t\in[0,T_0]}|X_t^\gamma-X_t^0|^2\right]
\le \mathbb{E}\!\left[\beta_{T_0}^\gamma\right]
+2\int_0^{T_0}\mathbb{E}\!\left[|X_s^\gamma-X_s^0|^2\right]\,ds
+2\sqrt{2}\sigma\left(\int_0^{T_0}\mathbb{E}\!\left[|X_s^\gamma-X_s^0|^4\right]\,ds\right)^{1/4}\exp\{T_0\}.
\]

Given $\gamma\le1$, $\mathbb{E}[\beta_\gamma(t)]\le C\gamma$ and according to (3.4),
\[
\int_0^{T_0}\mathbb{E}\!\left[|X_s^\gamma-X_s^0|^2\right]\,ds\le C\gamma,\qquad
\left(\int_0^{T_0}\mathbb{E}\!\left[|X_s^\gamma-X_s^0|^4\right]\,ds\right)^{1/4}\le C\gamma^{1/2}.
\]

With this, the rate $\mathcal{O}(\gamma^{1/2})$ leading to (3.12) follows. \hfill $\square$

\todo[inline]{5 end}

\section{Appendix.}

\noindent\textbf{Gronwall's lemma:}
Let $f,\alpha,g$ be three scalar functions, all defined on a certain interval $[0,T]$ and which are such that
\[
\forall\, t\in[0,T], \qquad f(t)\le g(t)+\int_0^t \alpha(s)f(s)\,ds.
\]
Then if $g$ is non-decreasing and $\alpha\ge0$, we have
\[
\forall\, t\in[0,T], \qquad f(t)\le g(t)\exp\left\{\int_0^t \alpha(s)\,ds\right\}.
\]

\noindent\textbf{Covariations/Crossvariation process and quadratic variation:}
Assume that $X$ and $Y$ are two real-valued continuous semi-martingales given by:
\[
X_t=X_0+M_t^X+A_t^X,\qquad Y_t=X_0+M_t^Y+A_t^Y,\ t\ge0,
\]
with $M^X$ and $M^Y$ being the respective martingale parts of $X$ and $Y$ and $A^X$ and $A^Y$ the finite variation parts. Then the covariation (crossvariation) process of $X$ and $Y$ is the stochastic process $\langle X,Y\rangle=(\langle X,Y\rangle_t)_{t\ge0}$ defined at all time $t$ by
\[
\langle X,Y\rangle_t=\limsup_{|t_{i+1}-t_i|\to0}\sum_{k=0}^{N-1}(X_{t_{k+1}}-X_{t_k})(Y_{t_{k+1}}-Y_{t_k}).
\]

The particular process $\langle X\rangle=\langle X,X\rangle$ corresponds to the quadratic variation process of $X$.

\noindent\textbf{It\^o's integration by parts:}
Let $(X_t;\,t\ge0)$ and $(Y_t;\,t\ge0)$ be two continuous semi-martingales. Then a.s., for all $t\ge0$,
\[
X_tY_t=X_0Y_0+\int_0^t X_s\,dY_s+\int_0^t Y_s\,dX_s+\langle X,Y\rangle_t.
\]

In the particular case where $X_t=\int_0^t \sigma_s\,dW_s$ and $Y_t=\int_0^t b_s\,ds$ with $\sigma$ and $b$ integrable $\mathcal{F}_t$-adapted processes, the above gives:
\[
X_tY_t=\int_0^t X_s b_s\,ds+\int_0^t Y_s\sigma_s\,dW_s,\qquad t\ge0.
\]

\noindent\textbf{It\^o's formula for general semi-martingale processes:}
Let $(X_t;\,t\ge0)$ be a $\mathbb{R}^d$-valued continuous semi-martingale with the decomposition:
\[
X_t=X_0+M_t+A_t,\ t\ge0,
\]
where $M$ is a martingale and $A$ a c\`adl\`ag process with finite variations. Then, for all scalar function
$f:[0,T]\times\mathbb{R}^d\to\mathbb{R}$, satisfying the following properties:
\begin{itemize}
\item for all $x$, $t\mapsto f(t,x)$ is differentiable and the partial derivative $(t,x)\mapsto \partial_t f(t,x)$ is continuous,
\item for all $t$, $x\mapsto f(t,x)$ is twice continuously differentiable on $\mathbb{R}^d$ and, for all $1\le i,j\le d$ $(t,x)\mapsto \partial_{x_i x_j}^2 f(t,x)$ is continuous,
\end{itemize}
it holds
\[
f(t,X_t)=f(0,X_0)+\int_0^t \partial_s f(s,X_s)\,ds+\sum_i\int_0^t \partial_{x_i}f(s,X_s)\,dX_s^i+\frac{1}{2}\sum_{i,j}\int_0^t \partial_{x_i x_j}^2 f(s,X_s)\,d\langle X^i,X^j\rangle_s.
\]

\noindent\textbf{Equivalent formulation:}
For $\nabla f$ the gradient of $f$, $\nabla_x^2f$ the Hessian of $f$ and $\langle X,X\rangle_t$ the $d\times d$ covariation matrix of $X$ (i.e. $\langle X,X\rangle_t^{(i,j)}=\langle X^i,X^j\rangle_t$)
\[
f(t,X_t)
\]
\[
=f(0,X_0)+\int_0^t \partial_s f(s,X_s)\,ds+\int_0^t \nabla_x f(s,X_s)\,dM_s+\int_0^t \nabla_x f(s,X_s)\,dA_s+\frac{1}{2}\int_0^t \mathrm{Trace}\bigl(\nabla_x^2 f(s,X_s)\,d\langle X,X\rangle_s\bigr).
\]
(this formulation particularly yields: under some general assumptions, the transformation $f(t,X_t)$ of a semi-martingale remains a semi-martingale.)

\todo[inline]{6 end}

\end{document}

